\documentclass[11pt]{article}
\usepackage[margin=1in]{geometry}
\usepackage{amsmath}
\usepackage{enumitem}

\title{Lesson Plan: Connecting Factors and Zeros}
\date{}

\begin{document}

\maketitle

\section*{I. Learning Target (Argue)}

Students will argue and explain the relationship between the zeros of a polynomial and its factored form.

\textbf{Students will be able to:}

\begin{enumerate}[label=\alph*.]
\item Articulate the relationship between zeros and the factored form of a polynomial.
\item Find the zeros of a polynomial given the factored form.
\item Explain the logical reasoning behind why a factor must equal zero when a product equals zero.
\end{enumerate}

\section*{Activity I: Exploration (20 minutes)}
\textbf{Teacher: Benjamin}

\subsection*{Goal}
Guide students to discover the connection between factors and zeros using graphs and factored expressions.

\subsection*{Steps}

\textbf{1. Identify Zeros (5 minutes)}

Students examine the graphs of the given polynomials.

Prompt students to:
\begin{itemize}
\item List all the zeros (x-intercepts) shown on each graph.
\end{itemize}

\textbf{2. Connect Zeros to Factors (10 minutes)}

Students match each zero with a factor in the polynomial.

Guiding prompts:
\begin{itemize}
\item Which factor becomes zero at each x-intercept?
\item What value of $x$ makes that factor equal to zero?
\end{itemize}

Students observe that each zero corresponds to a factor that becomes zero.

\textbf{3. Deepening Questions (3 minutes)}

Ask students:
\begin{itemize}
\item Are there any additional zeros outside the graphing window?
\item Why or why not?
\end{itemize}

Guide students to reason that a polynomial written as a product of factors becomes zero only when one factor equals zero.

\textbf{4. Formal Conclusion (2 minutes)}

State and record the conclusion:

\[
(x-a)(x-b)(x-c)\cdots = 0
\]

then

\[
x = a, b, c, \dots
\]

are the zeros of the polynomial.

Emphasize the \textbf{zero-factor condition}.

\section*{Activity II: Practice (25 minutes)}
\textbf{Teacher: Karina}

\subsection*{Problem Practice}

\textbf{Problems 1--4 (4 minutes)}

Students solve independently by applying the zero-factor condition.

Quickly review answers with the class.

\textbf{Problem 5 (5 minutes)}

Students compare the equation to the form discussed in Activity I.

Prompt students to identify differences such as:
\begin{itemize}
\item coefficients inside factors
\item factors not in the simple form $(x-a)$
\end{itemize}

Students adapt the method by setting each factor equal to zero.

\textbf{Problem 6 (5 minutes)}

Ask students:
\begin{itemize}
\item Is the polynomial written in factored form?
\end{itemize}

Students brainstorm possible strategies (e.g., factoring or quadratic methods).

Explain that this type of problem will be addressed later and move on.

\textbf{Problem 7 (5 minutes)}

Direct students’ attention to the factors:
\[
x \quad \text{and} \quad (3-x)
\]

Students apply the zero-factor condition again to determine the zeros.

\textbf{Problem 8 (6 minutes)}

Focus on the factor:
\[
x^2 - 4
\]

Guide students through:
\[
x^2 - 4 = 0
\]

\[
x^2 = 4
\]

\[
x = 2, -2
\]

Conclude that
\[
x^2 - 4 = (x-2)(x+2)
\]

Reinforce that factoring reveals the zeros.

\section*{Closing (Optional)}

Ask students:

\begin{quote}
Why does setting each factor equal to zero find the solutions?
\end{quote}

Students explain using the zero-factor condition.

\end{document}